\chapter{相关软件及技术简介}

本章介绍 CulCloud 平台开发过程中使用的主要技术，包括前端框架、后端服务、对象存储、任务队列、文件转换、PDF 渲染和大数据分析组件。

\section{前端开发技术}

\subsection{React 组件化开发}

\begin{figure}[H]
  \centering
  
\includegraphics[width=0.42\textwidth,height=3.2cm,keepaspectratio]{logos-pdf/logo-react.pdf}
  \caption{React 标识}
  \label{fig:chap02-react-logo}
\end{figure}

前端采用 React 框架构建模块化的单页面应用\cite{reactdocs}。系统将 Dashboard、MyFiles、ConvertCenter、PDFStudio 等业务视图抽象为独立组件，并设计了全局 DashboardContext 机制（见 file-cloud-frontend/src/context/useDashboard.tsx），统一维护视图流转与用户角色。在角色切换时，系统触发组件树重渲染，并在顶层容器动态绑定 theme-admin 样式，实现界面主题与视图权限的平滑过渡。

\subsection{TypeScript 类型约束}

\begin{figure}[H]
  \centering
  
\includegraphics[width=0.38\textwidth,height=3cm,keepaspectratio]{logos-pdf/logo-typescript.pdf}
  \caption{TypeScript 标识}
  \label{fig:chap02-typescript-logo}
\end{figure}

系统引入 TypeScript 进行静态类型约束，所有通信接口与状态协议均在 file-cloud-frontend/src/types.ts 中定义。通过对 FileItem、Task 等核心数据结构进行显式建模，在编译期捕获类型不匹配与属性缺失问题，保障文件预览、异步转换等多分支业务流转时的运行时稳定性。

\subsection{Vite 与界面构建}

\begin{figure}[H]
  \centering
  
\includegraphics[width=0.38\textwidth,height=3cm,keepaspectratio]{logos-pdf/logo-vite.pdf}
  \caption{Vite 标识}
  \label{fig:chap02-vite-logo}
\end{figure}

Vite 负责前端代码的开发调试与最终构建。利用其原生的 ESM 按需编译能力提升频繁交互界面（如 PDF 批注工作区）的响应开发效率。构建产物输出为优化后的静态资源，交由服务端托管部署。

\subsection{HTML5 Canvas 交互渲染}

\begin{figure}[H]
  \centering
  
\includegraphics[width=0.34\textwidth,height=2.7cm,keepaspectratio]{logos-pdf/logo-html5-canvas.pdf}
  \caption{HTML5 Canvas 标识}
  \label{fig:chap02-canvas-logo}
\end{figure}

HTML5 Canvas 是 PDF 工作区与管理驾驶舱的可视化基础\cite{mdncanvas}。在 PDF 批注中，系统在 Canvas 层上记录用户画笔、矩形和文本的相对坐标，并在导出时传给后端换算。同时，管理端 ECharts 默认采用 CanvasRenderer 绘制拓扑与统计图表，提供高效的可视化呈现。

\section{后端服务技术}

\subsection{FastAPI 接口服务}

\begin{figure}[H]
  \centering
  
\includegraphics[width=0.38\textwidth,height=3cm,keepaspectratio]{logos-pdf/logo-fastapi.pdf}
  \caption{FastAPI 标识}
  \label{fig:chap02-fastapi-logo}
\end{figure}

后端基于 FastAPI 构建高效的 REST API 路由\cite{fastapi}。利用 Pydantic 结构化校验前端提交的各类请求（如 PDF 页面重组请求，详见 backend/app/models/schemas.py），确保路由入口层的数据安全性；配合依赖注入机制，解耦鉴权、MinIO 客户端与任务状态调度逻辑。

\subsection{Flask 分析服务}

\begin{figure}[H]
  \centering
  
\includegraphics[width=0.34\textwidth,height=2.7cm,keepaspectratio]{logos-pdf/logo-flask.pdf}
  \caption{Flask 标识}
  \label{fig:chap02-flask-logo}
\end{figure}

Flask 框架用于承载独立的分析数据接口\cite{flaskdocs}。该分析服务与 FastAPI 解耦，专门用于读取 Spark 分布式聚合产物及服务遥测数据，并将其映射为管理驾驶舱所消费的历史指标接口，从而实现主业务流与分析遥测流在架构上的分层。

\subsection{Redis 与 Celery 任务队列}

\begin{figure}[H]
  \centering
  \begin{minipage}{0.38\textwidth}
    \centering
    
\includegraphics[height=3cm,keepaspectratio]{logos-pdf/logo-redis.pdf}
  \end{minipage}
  \hspace{0.08\textwidth}
  \begin{minipage}{0.38\textwidth}
    \centering
    
\includegraphics[height=3cm,keepaspectratio]{logos-pdf/logo-celery.pdf}
  \end{minipage}
  \caption{Redis 与 Celery 标识}
  \label{fig:chap02-redis-celery-logo}
\end{figure}

Redis 与 Celery 协同实现文件转换的异步任务队列。Redis 作为消息 Broker 并缓存近期任务的运行状态（含 queued, processing, success, failed，详见 backend/app/services/task\_queue.py）。通过将计算密集型的文档转换（如 LibreOffice 转换）从同步 HTTP 请求中解耦，前端可发起状态轮询展示进度，保证了核心接口的响应性能。

\subsection{MinIO 对象存储}

\begin{figure}[H]
  \centering
  
\includegraphics[width=0.42\textwidth,height=3.2cm,keepaspectratio]{logos-pdf/logo-minio.pdf}
  \caption{MinIO 标识}
  \label{fig:chap02-minio-logo}
\end{figure}

系统采用 MinIO 充当本地对象湖存储\cite{minio}。后端通过“对象名路径前缀”（详见 backend/app/routers/folders.py）模拟分层文件夹，重命名或移动时，级联拷贝前缀匹配的所有对象再批量清除源对象，在对象存储扁平架构上提供了高效的层级文件目录视图。

\section{文件处理技术}

\subsection{Office 与 PDF 转换}

\begin{figure}[H]
  \centering
  \begin{minipage}{0.38\textwidth}
    \centering
    
\includegraphics[height=3cm,keepaspectratio]{logos-pdf/logo-libreoffice.pdf}
  \end{minipage}
  \hspace{0.08\textwidth}
  \begin{minipage}{0.38\textwidth}
    \centering
    \includegraphics[height=3cm,keepaspectratio]{logos-pdf/logo-gotenberg.pdf}
  \end{minipage}
  \caption{LibreOffice 与 Gotenberg 标识}
  \label{fig:chap02-libreoffice-gotenberg-logo}
\end{figure}

系统集成 Gotenberg 转换服务处理异构文档格式\cite{gotenberg}。该服务容器化集成了 LibreOffice 和 PDF 引擎，后端通过 HTTP 接口提交 Office 转 PDF 任务，避免了直接在主业务进程中拉起本地进程的开销，结合 P0 格式白名单确保转换队列的可靠运行。

\subsection{CSV、Markdown 与文本渲染}

对于 CSV、Markdown 和 TXT 等轻量文本，系统采用本地轻量解析技术。后端直接提取文本并将其渲染为隔离的 HTML，前端通过预览容器直接呈现，无需触发繁重的文档转换流程。

\subsection{PDF 页面渲染与批注}

\begin{figure}[H]
  \centering
  
\includegraphics[width=0.38\textwidth,height=3cm,keepaspectratio]{logos-pdf/logo-pdfjs.pdf}
  \caption{PDF.js 标识}
  \label{fig:chap02-pdfjs-logo}
\end{figure}

对于 PDF 页面渲染与管理，系统以 PDF.js 的逻辑视口机制作为设计参照\cite{pdfjs}。PDFStudio 采用双层画布模型：底层展示后端 PyMuPDF 导出的高清页面图像，上层在 HTML5 Canvas 叠加手写及文本批注。前端记录相对视口的逻辑坐标与分辨率，导出时由后端进行坐标变换并合并写入新文件。

\section{运行与部署技术}

\subsection{PM2 进程管理}

\begin{figure}[H]
  \centering
  
\includegraphics[width=0.40\textwidth,height=3cm,keepaspectratio]{logos-pdf/logo-pm2.pdf}
  \caption{PM2 标识}
  \label{fig:chap02-pm2-logo}
\end{figure}

PM2 用于管理应用层的前端与分析模块服务\cite{pm2docs}。根目录下的 ecosystem.config.js 统一配置了开发与生产的服务进程。PM2 支持异常退出自动拉起与集中日志查阅，保障本地开发和现场验证时的服务连续性。

\subsection{Docker Compose 容器编排}

\begin{figure}[H]
  \centering
  
\includegraphics[width=0.40\textwidth,height=3cm,keepaspectratio]{logos-pdf/logo-docker-compose.pdf}
  \caption{Docker Compose 标识}
  \label{fig:chap02-docker-compose-logo}
\end{figure}

Docker Compose 用于编排系统底层的 Redis、MinIO、Gotenberg 及 API 镜像\cite{dockerdocscompose}。通过将持久化存储挂载、网络端口与多服务拓扑写入 docker-compose.yml 声明式配置，系统形成了“PM2 监管上层应用、Compose 监管基础设施”的混合部署架构。

\section{大数据分析技术}

\subsection{HDFS 与 Spark}

\begin{figure}[H]
  \centering
  \begin{minipage}{0.38\textwidth}
    \centering
    
\includegraphics[height=3cm,keepaspectratio]{logos-pdf/logo-apache-hadoop.pdf}
  \end{minipage}
  \hspace{0.08\textwidth}
  \begin{minipage}{0.38\textwidth}
    \centering
    
\includegraphics[height=3cm,keepaspectratio]{logos-pdf/logo-apache-spark.pdf}
  \end{minipage}
  \caption{Hadoop 与 Spark 标识}
  \label{fig:chap02-hadoop-spark-logo}
\end{figure}

HDFS 与 Spark 构成了本平台的大数据离线遥测分析管线\cite{shvachko2010hdfs,zaharia2016spark}。遥测流水日志进入存储后，Spark 计算引擎执行周期性的数据清洗、操作统计、错误率计算和格式流转矩阵分析，聚合出的指标存入静态 JSON，打通了“业务日志-数据清洗-指标分析-结果应用”的数据流。

\subsection{ECharts 可视化}

\begin{figure}[H]
  \centering
  
\includegraphics[width=0.38\textwidth,height=3cm,keepaspectratio]{logos-pdf/logo-apache-echarts.pdf}
  \caption{Apache ECharts 标识}
  \label{fig:chap02-echarts-logo}
\end{figure}

前端集成 ECharts 图表库进行多维度数据呈现。通过对实时事件流与 Spark 历史数据进行提取，使用折线图展示流量趋势、柱状图展示格式分布、雷达图展示质量监控、Sankey 拓扑图展示格式流转，提升了管理端驾驶舱的可视化运营密度。

\section{团队分工与协作机制}

为确保 CulCloud 云盘与大数据分析平台各项复杂技术的顺利落地与系统级集成，项目团队制定了高度工程化的团队分工与协作机制。各小组成员具体分工如下：

\textbf{曹磊（组长）}：全面负责系统的基础架构设计与核心后端开发。设计并实现了基于 FastAPI 的 REST API 接口体系、MinIO 级联目录与对象湖迁移逻辑、Redis 与 Celery 异步任务队列调度，以及 Gotenberg 文档转换引擎的容器化集成；同时，主导了基于 Apache Spark 的分布式离线分析管线，设计了管理驾驶舱的实时拓扑图与运行日志流面板，并编写了覆盖系统核心交互链路的 11 项 E2E 自动化测试用例，为代码集成合并提供了严格的质量保障。

\textbf{陈佳明（组员）}：负责用户端基础页面与文件管理前端交互开发。实现了“个人首页”与“我的文件”模块的前端渲染，包括多维度分类过滤器、路径面包屑导航、文件拖拽上传列表等 React 组件编写，并完成了全站的 CSS 样式适配与响应式布局调整，实现了用户端界面的平滑流转。

\textbf{陈晓楠（组员）}：负责 PDF Studio 处理面板前端视图开发。实现了 PDF 卡片拖拽重排与旋转/删除的状态管理，编写了 Canvas 轻量化手写及文本批注交互逻辑，并与组长协作调试了前端视口逻辑坐标到后端物理 PDF 页面坐标的换算与批注写回逻辑。

\textbf{程志阳（组员）}：负责格式转换控制台与分析接口辅助开发。实现了“转换中心”前端面板，编写了格式转换任务提交与基于定时器的状态轮询终止逻辑，同时协助组长开发了 Flask 离线分析微服务中的指标读取与遥测数据查询接口。

为了在系统验收与评审中展现系统工程的严密性，防止因代码细节引发追问，团队制定了高度工程化的协作与版本管理策略。首先，团队采用 Schema-First（接口先行）协作模式，开发初期通过 FastAPI 的 OpenAPI 文档，基于 Pydantic 明确定义了所有 API 请求与响应的 JSON 数据格式（Schema），形成了前后端的强契约关系，允许前端组员使用 Mock 数据并行开发视图，而组长在后端同步实现真实业务逻辑。其次，在版本管理上采用 Git Flow 分支管理规范，由组长曹磊维护主分支 v2-platform，组员分别在各自的特性分支（如 feat-myfiles、feat-pdfstudio）上独立进行提交，互不干扰。最终在分支合并时，组员必须在本地运行通过组长编写的 Playwright 自动化测试集，所有用例通过后方可由组长合并至主分支，通过严密的质量网关保障了代码的稳定集成。

\section{本章小结}

本章介绍了系统涉及的前端、后端、对象存储、任务队列、文件转换、PDF 渲染、运行部署和大数据分析技术。这些技术共同支撑 CulCloud 从单纯文件列表扩展为可上传、可预览、可转换、可编辑和可分析的资料处理平台。

从工程分层看，React、TypeScript、Vite 和 Canvas 负责构建浏览器侧的高交互界面，FastAPI 负责承载文件处理主链路，Flask 则作为轻量分析服务承担驾驶舱数据接口，二者形成“主业务服务”和“分析辅助服务”的边界。Redis 与 Celery 将耗时转换任务从同步请求中拆离，MinIO 负责对象化存储，Gotenberg、LibreOffice 和 PDF.js 分别补足 Office 转换、PDF 生成和页面渲染能力。PM2 与 Docker Compose 的组合则对应本地展示和容器化依赖两类运行场景，使系统既能快速预览，也能保持转换、缓存和对象存储等基础服务的可复现性。

这些技术并不是孤立堆叠，而是围绕文件流转闭环进行组合：用户上传资料后，后端写入对象存储并记录元数据；需要预览或转换时，任务队列调度转换服务生成结果；需要批注时，前端通过 PDF 页面图和 Canvas 批注层维护坐标关系；需要管理分析时，日志与遥测数据进入 Spark 聚合链路并由 ECharts 呈现。后续章节将在此基础上进一步说明需求建模、总体架构、核心模块实现和测试验证方法。
